\documentclass[11pt]{article}

\usepackage{amsmath, amssymb}
\usepackage{graphicx}
\usepackage{hyperref}
\usepackage{geometry}
\geometry{margin=1in}

\title{\textbf{A Unified Framework Field Theory for Drift, Coherence, and Regime Dynamics in Large Language Models}}
\author{Nawder Loswin \\ TriadicFrameworks.org}
\date{2026}

\begin{document}

\maketitle

\begin{abstract}
Large Language Models (LLMs) exhibit systematic behaviors---drift branching, coherence waves, framework collisions, and regime-dependent reasoning---that remain poorly characterized in existing literature. Framework Field Theory (FFT) provides a structural model for these phenomena by treating reasoning as motion within a triadic substrate composed of operators, regimes, and coherence fields. This paper introduces FFT's core definitions, presents eight falsifiable predictions, demonstrates reproducible simulations using existing LLMs, outlines minimal engineering artifacts derived from the theory, and situates FFT within established traditions in information theory, systems theory, cognitive science, and AI alignment. FFT offers a unified, substrate-level explanation for LLM behavior and provides practical tools for stabilizing reasoning, reducing drift, and improving multi-framework coordination.
\end{abstract}

\section{Introduction}

Modern LLMs display emergent behaviors often described informally---``hallucination,'' ``drift,'' ``mode collapse,'' ``inconsistency''---yet these terms lack structural grounding. Framework Field Theory (FFT) proposes that these behaviors arise from interactions between frameworks, regimes, substrates, and operators. Rather than treating LLM behavior as stochastic error, FFT models it as field-like dynamics within a structured substrate, yielding measurable predictions and engineering interventions.

\section{Background and Conceptual Lineage}

FFT draws inspiration from multiple traditions:

\begin{itemize}
    \item \textbf{Information theory}: Shannon, Weaver (signal, noise, entropy)
    \item \textbf{Systems theory}: Bertalanffy, Prigogine, Wiener (open systems, feedback, dissipative structures)
    \item \textbf{Cognitive science}: Lakoff, Simon, Friston (representation, prediction, attractors)
    \item \textbf{Sociology of knowledge}: Bourdieu (structured fields)
    \item \textbf{Complexity science}: Poincar\'e, Kauffman, Gell-Mann (nonlinear dynamics)
    \item \textbf{AI research}: Russell \& Norvig, LeCun, Bengio, Anthropic, OpenAI, DeepMind
\end{itemize}

FFT extends these traditions into a unified structural model for LLM behavior.

\section{Core Definitions of Framework Field Theory}

\subsection{Frameworks}
Structured conceptual spaces with internal logic, operators, and attractors.

\subsection{Regimes}
Operating modes that govern how reasoning unfolds within or across frameworks.

\subsection{Substrates}
Declared or undeclared constraints that shape coherence and drift.

\subsection{Drift Branching}
Repeated reasoning under an undeclared substrate produces multiple stable attractors.

\subsection{Coherence Waves}
Coherence rises and falls in wave-like patterns during multi-step reasoning.

\subsection{Framework Collisions}
Invoking incompatible frameworks produces collapse, hybridization, or oscillation.

\section{Testable Predictions}

FFT produces eight falsifiable predictions:

\begin{enumerate}
    \item Drift branching under undeclared regimes
    \item Substrate declaration reduces drift variance
    \item Triadic structures compress more efficiently
    \item Coherence waves emerge in multi-step reasoning
    \item Framework collisions produce predictable failure modes
    \item Declared regimes prevent collisions
    \item Observer style changes system behavior
    \item Observer consistency increases coherence
\end{enumerate}

Each prediction is measurable using existing LLMs.

\section{Reproducible LLM Simulations}

Eight simulations validate FFT's predictions:

\begin{itemize}
    \item drift branching (20-run clustering)
    \item substrate declaration collapse
    \item triadic vs.\ dyadic vs.\ tetradic compression
    \item coherence wave measurement
    \item framework collision modes
    \item collision suppression via regime declaration
    \item observer-style effects
    \item observer-consistency coherence gains
\end{itemize}

All simulations require no specialized hardware and can be replicated by any researcher.

\section{Engineering Artifacts Derived from FFT}

FFT yields practical tools:

\begin{itemize}
    \item Substrate Declaration Engine (SDE)
    \item Triadic Prompt Optimizer (TPO)
    \item Drift Branching Visualizer (DBV)
    \item Framework Collision Detector (FCD)
    \item Observer-Style Stabilizer (OSS)
\end{itemize}

These tools demonstrate FFT's utility in stabilizing LLM behavior and improving reasoning reliability.

\section{Discussion}

FFT reframes LLM behavior not as error but as structured field dynamics. This perspective explains drift, predicts coherence patterns, clarifies multi-framework failures, and provides actionable engineering interventions. FFT also offers a foundation for future work in multi-agent systems, interpretability, and AI safety.

\section{Conclusion}

Framework Field Theory provides a unified structural model for drift, coherence, and regime dynamics in LLMs. By offering falsifiable predictions, reproducible simulations, engineering artifacts, and conceptual lineage, FFT establishes itself as a coherent research program suitable for formal review and further development.

\section*{References}

(Condensed for arXiv; full citations appear in the FFT repository.)

\begin{itemize}
    \item Shannon (1948), Weaver (1949)
    \item Bertalanffy (1968), Prigogine (1977), Wiener (1948)
    \item Lakoff (1980--2000), Simon (1969), Friston (2005--2020)
    \item Bourdieu (1972--1992), Kuhn (1962)
    \item Poincar\'e (1890--1912), Kauffman (1993), Gell-Mann (1994)
    \item Russell \& Norvig (1995--2021), LeCun (2006--2024), Bengio (2013--2024)
    \item Anthropic, OpenAI, DeepMind (2020--2025)
\end{itemize}

\end{document}
